\section{A global Frobenius fills the degree-twenty-nine deficit}
\label{sec:global-frobenius-29}

Put \(\zeta=\zeta_{29}\), \(O=\mathbf Z_2[\zeta]\), and let
\(s(\zeta)=\zeta^2\). Since two has order twenty-eight modulo
\(29\), the ring \(O\) is unramified of degree twenty-eight.
Write \(c=s^{14}\), and put
\[
 U=1+2O,\qquad T=\{u\in U:u\,c(u)=1\},\qquad P=(s-1)T.
\]
The explicit CM-type operator is
\[
 Q=-s+s^3-s^4+s^5-s^7+s^{10}-s^{11}.
\]
The preceding integral calculation gives
\[
 (s-1)\Phi^*=(c-1)Q,\qquad Q(1)=-1,\qquad [T:QT]=8.
\]
Let \(\mathcal R_\Pi,\mathcal V_\Pi\) denote the established simultaneous kernels
for the primes satisfying the integral unit criterion. In particular
\(29\notin\Pi\), \(G=\mathcal R_\Pi H\), and \(\mathcal R_\Pi/\mathcal V_\Pi\) is
pro-odd. The compact-homology scalar in degree twenty-nine is
\(\mathfrak h_{29}\).

\begin{theorem}\label{thm:global-frobenius-29}
All four images are the whole target:
\[
 \mathfrak h_{29}(\mathcal R_\Pi)=\mathfrak h_{29}(\mathcal R_\Pi\cap H)
 =\mathfrak h_{29}(\mathcal V_\Pi)=\mathfrak h_{29}(\mathcal V_\Pi\cap H)=P.
\]
In this inert case \(P\simeq\mathbf Z_2^{14}\).
For
\[
 \mathcal R_{\Pi,29}=\mathcal R_\Pi\cap\ker\mathfrak h_{29},\qquad
 \mathcal V_{\Pi,29}=\mathcal V_\Pi\cap\ker\mathfrak h_{29},
\]
one has \(G=\mathcal R_{\Pi,29}H\). Both new kernels have infinite
index, with exact quotient \(P\), already on their arithmetic
intersections. They preserve respectively the full relative quotient
\(B\) and its maximal pro-two quotient. The original invariant
quadratic-character restriction kernel remains exactly \(\mathcal S\), and
its image on either new kernel is exactly the invariant quadratic
characters killing the arithmetic intersection, canonically
\(\operatorname{Hom}_{\rm cts}(B,\mathbf F_2)\).
\end{theorem}

\subsection{The exact missing finite quotient}

The torus coinvariant calculation gives \(T/P\simeq C_2^2\),
and \(s\) acts trivially there. Thus \(Q\) induces multiplication
by \(Q(1)=-1\), an automorphism. Inclusion induces an isomorphism
\begin{equation}\label{eq:g29-cokernel-identification}
                         P/QP\xrightarrow{\sim}T/QT.
\end{equation}
Indeed, if \(Qt\in P\), invertibility on \(T/P\) forces
\(t\in P\); this proves injectivity. For any \(t\in T\),
choose \(u\) with \(Q[u]=[t]\) in \(T/P\), and then
\(t/(Qu)\in P\), proving surjectivity.

Set \(T(4)=T\cap(1+4O)\). The first principal-unit layer is
\[
 \theta:T\longrightarrow \mathbf F_{2^{14}},\qquad
             t\longmapsto\frac{t-1}{2}\pmod2,
 \qquad \ker\theta=T(4).
\]
The norm condition puts the image in the fixed field of \(c\).
This map is onto: lift any element of that fixed field to \(a\in O\).
The norm of \(1+2a\) lies in \(1+4O^c\). The norm on
\(1+4O\) is onto \(1+4O^c\), by logarithm and surjectivity
of the trace in the unramified quadratic extension. Correcting by
an element of \(1+4O\) therefore produces a norm-one lift with
the specified first layer.

The fourteen elements
\[
                   \beta_j=\zeta^j+\zeta^{-j},\qquad 1\le j\le14,
\]
are an \(\mathbf F_2\)-basis of this fixed field. They are independent
because \(\zeta,\ldots,\zeta^{28}\) are a basis of
\(O/2O\), and these sums have disjoint pairs of supports.
Frobenius cycles the normal basis \(\beta_{2^k}\), \(0\le k<14\),
where indices are read modulo twenty-nine and up to sign. Thus, as an
\(\mathbf F_2[s]\)-module,
\[
                    \mathbf F_{2^{14}}\simeq
                         \mathbf F_2[s]/(s^{14}+1).
\]
The exact polynomial calculation is
\[
 \gcd(\overline Q,s^{14}+1)=f(s),\qquad f(s)=s^3+s^2+1.
\]
Consequently
\(\mathbf F_{2^{14}}/\overline Q\mathbf F_{2^{14}}
\simeq\mathbf F_2[s]/(f)\), of order eight.
The induced map from \(T/QT\) is surjective; since its domain also
has order eight, it is an isomorphism. Together with
\eqref{eq:g29-cokernel-identification}, this identifies the missing
quotient exactly as
\begin{equation}\label{eq:g29-layer-certificate}
                         P/QP\simeq\mathbf F_2[s]/(s^3+s^2+1).
\end{equation}
In particular it has exponent two. The cubic polynomial is irreducible,
so the \(s\)-orbit of every nonzero element spans the quotient.

\subsection{The arithmetic word allows arbitrary CM Frobenius}

Choose \(a\in I_{29}\) with \(\lambda(a)\bmod29=2\).
For any \(b\in G_{\mathbf Q(\zeta)}\), put
\[
                       z=[a,b],\qquad z_2=z^{e_2},\qquad w=[a,z_2],
\]
where \(e_2\in\widehat{\mathbf Z}\) is one at two and zero at
all other primes, with exponentiation in the procyclic closure.
The element \(b\) is not required to be an inertia element.

Both \(z\) and \(z_2\) have full multiplier one. The
coefficient-two Kummer cocycle of two vanishes on \(a\in I_{29}\),
and \(\lambda_2(a)=1\). The cocycle commutator formula therefore
gives \(\rho_{2,2}^I(z)=0\). Primary projection kills every other
coefficient, so \(z_2\in N\cap H\). Every rational Kummer
coordinate of \(z_2\) is supported at two, and hence
\[
                 \rho_q^I(w)=(\lambda(a)-1)\rho_q^I(z_2)=0
                         \quad\text{for all rational primes }q.
\]
Thus \(w\in D\cap H\).

At coefficient two the based action of \(a\) is trivial, by
unramifiedness outside two, so \(\Phi_2(z_2)=\Phi_2(w)=1\).
For every odd prime \(\ell\), the normalized image of a
multiplier-one element lies in the pro-\(\ell\) kernel of
\(\operatorname{Aut}(\widehat F_2^{(\ell)})\to
\operatorname{GL}_2(\mathbf F_\ell)\). The image of the pro-two
element \(z_2\) in that group is therefore trivial.
For completeness, the pro-\(\ell\) assertion follows on finite
characteristic quotients from the following elementary fact. The
Frattini kernel in the automorphism group of a finite \(\ell\)-group
acts freely on the set of generating tuples lifting a fixed basis of
the Frattini quotient. That nonempty set has \(\ell\)-power cardinality,
so the automorphism kernel is an \(\ell\)-group. Passing to
characteristic quotients proves the assertion.
Hence all primary based word actions kill \(w\).

The element \(a\) fixes both elliptic CM fields
\(\mathbf Q(i)\) and \(\mathbf Q(\sqrt{-3})\), as does \(z_2\).
Their elliptic Tate actions therefore commute. The corrected harmonic
frame of \(w\) is one, and so is the corrected cubic frame, since
its quartic Kummer correction of three vanishes on \(D\).
The cubic radical--CM defect vanishes arithmetically. Finally
\[
                 \chi(w)=\chi(a z_2a^{-1})-\chi(z_2)=0
                              \quad(\chi\in X_2).
\]
All the quadratic coordinates defining \(V_{\rm quad}\) vanish.
Thus \(w\in V_{\rm quad}\cap H\).

For an odd prime \(q\ne29\), both \(a\) and \(z_2\) have
trivial degree-\(q\) deck exponent. Their compact-homology actions
lie in the commutative deck centralizer. Consequently
\(\mathfrak h_q(w)=1\) for every such \(q\), and
\begin{equation}\label{eq:g29-word-membership}
                               w\in\mathcal V_\Pi\cap H.
\end{equation}
This verifies all previous kernels for arbitrary \(b\) fixing
\(\mathbf Q(\zeta)\).

If \(\alpha=\rho_{J_{29},2}(b)\), the Jacobian value of the first
commutator is \((s-1)\alpha\). The based-cover/Jacobian comparison
may have a deck \(C_{29}\) discrepancy; it vanishes after pro-two
projection. The full cover representation then gives exactly
\begin{equation}\label{eq:g29-word-value}
                  \mathfrak h_{29}(w)=\bigl((s-1)^2\alpha\bigr)^{e_2}.
\end{equation}

\subsection{The Frobenius at fifty-nine and its exact normalization}

The Fermat curve \(u^{29}+v^{29}=1\) maps to \(C_{29}\) by
\[
                             (u,v)\longmapsto(t=u^{29},w=-uv).
\]
It is the quotient by \((u,v)\mapsto(\zeta u,\zeta^{-1}v)\).
At coefficient two, averaging over this order-twenty-nine subgroup
identifies compact homology of the quotient with its invariants,
integrally. The surviving Fermat character components are exactly
\((a,a)\), \(a\ne0\). The degree-twenty-nine deck action on the
quotient is induced by \((u,v)\mapsto(\zeta u,v)\), so the
\((1,1)\) component has deck eigenvalue \(\zeta\).

By \cite[Theorem 8.1 and Remark 8.3, pp.~638--639]{Otsubo2016},
arithmetic Frobenius acts on this homology component, equivalently
on the Jacobian Tate module, by
\[
 j_{29}^{1,1}(\mathfrak l)
       =-\sum_{x\in\mathbf F_{59}\setminus\{0,1\}}
                             \chi_{\mathfrak l}(x)\chi_{\mathfrak l}(1-x).
\]
The cited remark explicitly distinguishes geometric Frobenius,
cohomology, and homology; the scalar here is the displayed Jacobi
sum, not its inverse. Choose \(\mathfrak l\mid59\) with
\(\zeta\equiv4\pmod{\mathfrak l}\). Since two is primitive
modulo fifty-nine, the residue character satisfies
\(\chi_{\mathfrak l}(2)=\zeta\). Choose \(b\) lifting arithmetic
Frobenius at this prime. It fixes \(\mathbf Q(\zeta)\), and its
coefficient-two scalar is exactly
\begin{equation}\label{eq:g29-jacobi}
 \alpha=\zeta^{27}
       +2\sum_{j\in\{4,5,9,13,14,24,26\}}\zeta^j
       -2\sum_{j\in\{7,10,17,19,22,23,25\}}\zeta^j.
\end{equation}
For a direct finite verification, the counts of the exponents
\(\log_2x+\log_2(1-x)\pmod{29}\), in order from zero to
twenty-eight, are
\[
\begin{gathered}
(2,2,2,2,0,0,2,4,2,0,4,2,2,0,0,\\
2,2,4,2,4,2,2,4,4,0,4,0,1,2).
\end{gathered}
\]
Subtracting the final coefficient with the relation
\(\sum_{j=0}^{28}\zeta^j=0\), and taking the minus sign,
gives \eqref{eq:g29-jacobi}. Exact multiplication also gives
\(\alpha\bar\alpha=59\).

Set \(\delta=\alpha\zeta^2\). Then \(\delta\in1+2O\), and
\[
              (s-1)^2\alpha=\zeta^{-2}(s-1)^2\delta.
\]
The first factor has odd order, while the second is pro-two.
Equation \eqref{eq:g29-word-value} consequently becomes the exact
identity \(\mathfrak h_{29}(w)=(s-1)^2\delta\).
Writing \(\delta=1+2r\), its first layer is
\[
 \theta(\mathfrak h_{29}(w))=r+s^2r
                   =\beta_2+\beta_{11}+\beta_{12}+\beta_{13}.
\]
The last equality follows by substituting
\eqref{eq:g29-jacobi} and reducing modulo four. In the cyclic
basis \(\beta_{2^k}\), it is the polynomial
\[
                   A(s)=s+s^4+s^7+s^{11}.
\]
Its remainder modulo \(s^3+s^2+1\) is \(s+1\), which is
nonzero. Thus this arithmetic Frobenius word has a
nonzero image in \(P/QP\). Its first three conjugates by \(a\)
have classes
\[
                        s+1,\qquad s^2+s,\qquad 1,
\]
which form a basis of that three-dimensional quotient.

\subsection{Saturation and the new relative kernel}

The local product-inertia construction already realizes \(QP\)
on \(\mathcal V_\Pi\cap H\): its exact value is
\((s-1)^2\Phi^*(U)=(s-1)QT=QP\), and its other
prime-degree coordinates are one. The Frobenius word just
constructed and its conjugates supply every class in \(P/QP\).
Therefore \(\mathfrak h_{29}(\mathcal V_\Pi\cap H)=P\).
The full-group upper bound \(\mathfrak h_{29}(\mathcal R_\Pi)\subset P\)
proves all four image equalities in the theorem.

Arithmetic correction in this common image proves
\(G=\mathcal R_{\Pi,29}H\), and ordinary invariant-character gluing
preserves the full relative quotient \(B\). The quotient
\(\mathcal R_{\Pi,29}/\mathcal V_{\Pi,29}\) embeds in the old pro-odd quotient;
central coprime extension therefore preserves \(B_{\max2}\).
Every arithmetic element used to fill \(P\) is killed by every
original invariant quadratic character. Thus an original character
vanishing on a new kernel, and hence factoring through \(P\)
on the old subgroup, already vanishes on that old subgroup.
Its original restriction kernel is exactly \(\mathcal S\).
The complete quadratic splitting, gluing on \(\mathcal R_{\Pi,29}\), and
pro-odd extension for \(\mathcal V_{\Pi,29}\) give the stated complete
restricted-image description. Finally \(P\) is torsion-free of
rank fourteen in this single inert factor, so both reductions are
strict and have infinite index.
